% =====================================================================
%  Georg Brandes, HANS BRØCHNER (1886).
%  In: Georg Brandes, Essays. Danske Personligheder [Essays: Danish
%  Personalities] (Copenhagen: Gyldendal, 1889), pp. 230--255.
%
%  TRANSLATION (English).  Translated 2026-10-27 from transcription.tex in
%  this directory, as rebuilt that day (ebook-method/hans-brochner/; its
%  header records how the text was established).  Method:
%  ../../../TRANSLATION-PLAYBOOK.md.  \opage{N} marks where printed page N
%  begins, at the same joint as in transcription.tex (between the English
%  words nearest a word the print divides across the turn).
%
%  ---- REGISTER ----
%  Readable modern English that keeps Brandes's periods where English will
%  carry them.  The essay is a memoir-portrait, warm and exact; nothing
%  softened.
%
%  ---- CONVENTIONS ----
%  Italics in the print: emphasis -> \emph; titles of works -> \textit, the
%  Danish title kept with an English gloss in brackets.  Guillemets «...»
%  -> ``...''.  Danish verse translated in the text, the Danish in a comment
%  below it.  Videnskab = science; Docent = docent; Aand = mind or spirit
%  as the context asks (Universets Aand = the spirit of the universe).
%  The print's two misprints (p. 247 "i i", p. 252 a semicolon) are noted
%  in comments where they fall; the English translates the sense.
% =====================================================================
\documentclass[12pt,a4paper,openany]{book}
\usepackage[utf8]{inputenc}
\usepackage[T1]{fontenc}
\usepackage{libertinus}
\usepackage{libertinust1math}
\usepackage{textalpha}
\usepackage{graphicx}
\usepackage[english]{babel}
\usepackage[protrusion=true,expansion=true]{microtype}
\usepackage{setspace}
\usepackage[margin=1.3in]{geometry}
\usepackage{enumitem}
\usepackage{fancyhdr}
\usepackage{hyperref}
\hypersetup{pdftitle={Hans Brøchner (1886)}, pdfauthor={Georg Brandes}, hidelinks}
\setlength{\headheight}{15pt}
\pagestyle{fancy}
\fancyhf{}
\fancyhead[LE,RO]{\thepage}
\fancyhead[C]{\textit{Georg Brandes: Hans Brøchner (1886)}}
\setstretch{1.3}
\usepackage{marginnote}
\newcommand{\opage}[1]{\marginnote{\footnotesize\textit{[#1]}}}
% the essay's head as printed: title, date, a short rule (as in transcription.tex)
\newcommand{\essaytitle}[2]{\begin{center}{\large #1}\\[4pt]#2\\[4pt]\rule{1.5cm}{0.4pt}\end{center}}
\newcommand{\sectionnum}[1]{\begin{center}#1\end{center}}
% the double rule printed at the end of the essay
\newcommand{\closingrule}{\begin{center}\rule{3cm}{0.4pt}\\[-0.8\baselineskip]\rule{3cm}{0.4pt}\end{center}}

\begin{document}
\thispagestyle{empty}
\begin{center}
\vspace*{2cm}
GEORG BRANDES\\[2ex]
{\LARGE HANS BRØCHNER}\\[1ex]
(1886)\\[4ex]
from \textit{Essays. Danske Personligheder} [Essays: Danish Personalities],\\
pp.~230--255\\[\fill]
\end{center}
\vfill
\begin{center}
Copenhagen: Gyldendalske Boghandels Forlag (F.\ Hegel \& Søn), 1889\\[3ex]
{\small Translated from the Danish}
\end{center}
\clearpage
\noindent\textit{Contents} (p.~[v]) of the volume, the entry for this essay: H. Brøchner \dotfill\ 230

\vspace{2em}
% ===== TEXT (pp. 230--255) =====
% --- p. 230 ---
\opage{230}\essaytitle{HANS BRØCHNER.}{(1886.)}

\sectionnum{I.}

What one seeks from the teachers at a university at nineteen, with the first
years of study behind one, is not so much knowledge as the impression of a
strong, mature personality. One desires instruction less than the guidance
that can be given to the young by a mind well oriented in more than the world
of books. One wishes for a teacher who is also a character.

In 1861 there were a number of young men who glowed for philosophy; nearly all
of them believed in it. R. Nielsen was sought out by not a few. He always showed
interest and goodwill to those who came with questions. But Brøchner expressly
offered, in the lecture catalogue, guidance to anyone who wished to study
philosophy. That decided the matter for
% --- p. 231 ---
\opage{231}some to whom a coherent course of inquiry was a serious matter. And soon it
was not thirst for knowledge alone that drew us to Brøchner. His personality
drew us.

He was 41 years old when we first saw him, heard a series of lectures from
him, and were captivated by his elegant and austere figure, the classical
profile of his southern head, the great earnest gaze under the arched,
somewhat weary lids. Though somewhat worn by the years, he was the noblest
male beauty we had ever seen. He was tall and slender, lean and spare. The
face was narrow and somewhat hollow, the rich, soft hair blue-black in its
smoothness. The large brown eyes had an expression as though they had wept
much, as though their gaze had been softened by sorrow. Above them the
forehead rose steep and white, as if in two storeys. The Spanish beard did its
part to give the impression of something foreign. He did not look like a
modern scholar, rather like a thinker from the Italy of the Renaissance, and
was indeed, by his own account, of Italian descent.

He looked like a man for whom good living did not exist; he plainly made it a
point of honour, in true Spinozist fashion, not to have striven, nor to strive,
for earthly well-being. He lived in
% --- p. 232 ---
\opage{232}a philosophical poverty; for even when he was married his salary as docent
was only 2,400 kroner. His dress bore the stamp of strict frugality; it was
simple and dark, never a spot. His waistcoat was usually buttoned, against the
fashion of the day, right up to the neck. He wore the same clothes day in, day
out, and they looked the same; he seemed to wear nothing out. His appearance
recalled Welhaven's verse on the Polish émigré:

\begin{verse}\small
His face was faded, his dress was coarse, \\
His linen was clean and his hands were fine.
\end{verse}
% Danish: Hans Ansigt var falmet, hans Dragt var grov, /
%         Hans Linned var rent og hans Hænder var smukke.

It was through the eye, not through the ear, that one felt drawn to him from
the first. For his lecturing was anxiously correct, fastidiously groomed in
its language, abstract in its train of thought, and answered no question
otherwise than by a detailed account of the history of the question and of the
answers given to it. It moved forward monotonously, in measured sentences, and
had almost no character of speech, though it was delivered freely enough.
Brøchner did not have the gift of speech, which is the gift of giving each
single listener the impression that the speech is addressed to him in
particular. His words often went over our heads; they instructed us without
striking home. The lecture was consistent and firm in
% --- p. 233 ---
\opage{233}thought, but it sounded as though it were read aloud from an inner book; it
did not \emph{live}.

The speaker, on the other hand, lived for us. It was a personality who stood
there, as distinguished as he was interesting.

He awakened sympathy; for he looked as though he had once been all passion,
while now he was all thought. The sorrowful earnestness of his gaze drew one
in, while at the same time something embarrassed and shy that as a rule lay
over the speaker's bearing kept one at a distance and raised the question
whether one could come close to this man as a human being.

And he commanded respect. The little we young people knew of his fortunes
revealed an uncommon independence and strength of character. When, quite a
young man, he presented himself for the theological examination, he had shown
the openness and the courage to profess his view of the dogmatics of
Christianity as mythical. In consequence he had been forbidden to take that
examination. He had then thrown himself into the study of philosophy, had
translated Strauss, had relied on Feuerbach in a treatise on \textit{Daaben}
[Baptism] directed against Martensen, and was so hated as a dangerous man that
Christian the Eighth had with his own hand struck his name from the list of
applicants for a travel
% --- p. 234 ---
\opage{234}grant. Nonetheless he had (with the help of a brother) lived much abroad,
stayed long in Germany, been twice in Italy, and only after 1848 had he again
regarded Denmark as a country that had room for him.

The contrast in which his public appearances stood to Rasmus Nielsen's so much
livelier, but so far more virtuoso, lecturing struck everyone. It was a
contrast that must have corresponded fairly closely to the one that
Jouffroy's lectures at the Sorbonne in Paris in the thirties presented to
Cousin's --- if the following French portrait of the philosopher Jouffroy,
which we translate with some omissions but without adding anything whatever,
was otherwise true to reality:

``One entered Jouffroy's lecture hall and found there a lean man, somewhat
stooped, like all consumptives, with deep-set eyes, a thoughtful and
melancholy physiognomy, and an expression over his whole person of weariness,
nobility and resignation. He was proud, reserved, often silent, neither
forthcoming nor familiar with his students. While Cousin, when set in motion
by the storm of enthusiasm and the overflowing abundance of animal life,
talked, communicated himself, poured himself out,
% --- p. 235 ---
\opage{235}explained and persuaded with gestures and the whole oratorical apparatus, no
matter whether in a public garden or in his room, no matter where and no
matter for whom --- Jouffroy was cold and restrained even in the lecture
chair. He looked as though he did not know that anyone was there.

His gestures were sparing, his body motionless; one would have thought he was
reading an invisible book, occupied solely with understanding it, and wished
only to convince himself; he was thinking aloud. No brilliant words, no daring
sentences, no effort to amuse, rouse or move; on the contrary, long
introductions full of petty divisions and subdivisions, a circumstantial and
unceasing examination of the preliminary questions.

When at last he had penetrated into his subject, no decisive or hasty
proposition, endless repetitions. When his disciples read through their notes,
they found that they had written down the same idea three or four times. And
yet from the first day every attentive mind bowed beneath his mind. One
discovered at once in him a hidden hearth where an inextinguishable fire
burned, more violent and more powerful than Cousin's dazzling illumination.
The less he gave himself
% --- p. 236 ---
\opage{236}away, the more one felt how he burned. These redoubled repetitions showed,
in their stubborn attempts, a mind insatiable for the exact expression, unable
to rest content with thoughts merely sketched, invincible in the face of
effort, persevering in the course of its thoughts until it had reached
complete clarity, and at last well acquainted with the narrowness of the way
and with the streams of error that flow on both sides of it. One felt reverence
and trust, and when a tremor in the voice, or once in a while an image,
signalled an important truth he had found, one saw in this faint sign more
emotion and more eloquence than in his rival's brilliant dithyrambs.''

It would hardly be possible to give a more apt portrait of the merits of
Brøchner's teaching if one set out expressly to do so.

And critical as he was above all by disposition, he shone especially as a
polemicist. His offensive war against the theologians, or against Nielsen's
masked-theological philosophy, was a model of superior tactics. He would
usually begin by setting out with the greatest precision the view he meant to
contest, then present point by point the proofs of its untenability, then
adduce everything
% --- p. 237 ---
\opage{237}that could conceivably be urged from the opponent's side against those
proofs, and explain why this supposed defence would count for nothing; then he
proved that no defence other than the one adduced could be made at all.

When this was done, he dissolved his opponent's connection of concepts, made
him concessions, supposed \textit{per impossibile} that one of the overthrown
main points was a firm castle, and showed that even so the most decisive
categories contained contradiction upon contradiction. Step by step he went
round his opponent's edifice of doctrine, and, well acquainted with its manner
of fighting, stopped up every hole through which the opponent might try to
escape, until he had him securely caught. Then he would usually end by
recapitulating all the contradictions he had shown, and at last closed the
critique, held in the most serious tone, with an irony that one could in all
truth, and without abuse of the word, call terrible.

\sectionnum{II.}

A legend went about among the young students that Brøchner, in his young days
down in Italy, had been engaged to a very distinguished young lady, a princess
or a countess, but that the match was said to have foundered on
% --- p. 238 ---
\opage{238}the difference of religion. It gave him a romantic nimbus even now, when he
was engaged to a young Danish girl from the provinces, of middle-class
origin.

In 1862 he married, and lived in a happy marriage until 1866, when his young
wife died after giving birth to her second child. The two had the best of
relations, despite a difference in age of a full twenty years, and despite the
domain of abstract philosophy, inaccessible to a young woman, round which
Brøchner's thoughts, and as a rule his talk too, revolved. His marriage was an
idyllic but short episode.

Through a much longer period of his life the memory of that Italian woman,
Costanza T., whom he had known and loved in his youth, stood like a star over
his daily existence. She had become the dream of his life after the two had
been parted in their youth. A very beautiful daguerreotype of her as a young
woman always hung in his study. He spoke not infrequently of the original, as
of a young lady he had once known, without however going more closely into
their relationship.

In 1866 I came to know her personally. She was living then, married to a
countryman of hers, Count Oreste B., in Paris, but really
% --- p. 239 ---
\opage{239}outside Parisian life, in the company only of the Italian, Spanish and South
American colony.

Her beauty, but above all her apparent youth, surprised me. I had expected to
see an elderly lady, but she seemed no more than some twenty-six or
twenty-seven years old, and her figure was as youthful as her face.

One could hardly have had a better recommendation to her than Brøchner's. She
could not hear enough about his life and doings. When in our first
conversation she asked me whether Brøchner had grown old, and I answered that
his hair was still black, she added: And what a black it is! As an
intellectual individuality Brøchner did not really exist for her; she had not
penetrated into his nature at all, and could not even have made the attempt.
She was a young child when he came to know her, and even now, a grown and very
delicately feeling woman, she was at bottom childlike. She was unlettered, as
Italian women are, extremely natural and unaffected, but reserved none the
less, always in a mild and good humour against a background of a certain
melancholy. Not exactly happy, but content. She was calm, and by contrast with
her merry, sparkling, worldly and more brilliant, but less beautiful
% --- p. 240 ---
\opage{240}sister, in whose house she was staying, she seemed calmer and quieter still.
For several months at that time I saw her almost daily; in 1870 I stayed in
her house in the village of Alpignano near Turin, and in 1871 met her again in
Rome. She had the deep amiability that is never darkened by a mood; there was
the same gentle harmony in her soul as lay over her features.

Costanza was a refined fruit of the luxuriant soil of her native land. She
resembled the noble kind of grape that in Italy is called the strawberry
grape, which unites the rich sweetness of the grape's juice with the naive
aroma of the wild strawberry. She had been badly used in her relation to
Brøchner. Her mother, who was as vain as she was bigoted, had, after
Brøchner's departure, by constant pressure at home and through the influence
of the father confessor, had the engagement to the foreign heretic broken off.
Many years later she married a freethinker, as vehement as only Italy breeds
them, far more violent in his pronouncements than Brøchner ever was, and
became attached to him with all her heart, without however joining in his
convictions. Her husband was one of the most sensible, most upright and most
gifted Italians I
% --- p. 241 ---
\opage{241}have known, a hater and despiser of all sham. He valued Costanza as she
deserved.

She constantly nursed the thought of seeing Brøchner again, and many were the
plans made for journeys for him to that end. But the thought of a reunion did
not seem to tempt him. He never entered into it, and seemed rather to shun it.
It was hardly only the difficulty that long journeys would have meant for him
that held him back. Neither the lack of means for them nor his poor health
seemed decisive. Probably he shunned the emotion as painful. Or he had grown so
used to a life in the memory of his youth and in the dream-relation to
Costanza that he feared the disturbing breath of reality. He had to keep her
image at the distance of remoteness.

For her, as a woman and a southerner, this romanticism did not exist. She had
once loved him; she was still fond of him with all her heart. It would have
been a joy to her to see him again, to nurse him, now that he was so gravely
ill. She was very anxious that he should never form an unfavourable or false
opinion of her. When he married, she broke off the correspondence with him,
for fear that his wife might find something strange
% --- p. 242 ---
\opage{242}in it. But then, when the New Year came, she seized the occasion to send him
her good wishes. Later the correspondence was taken up again, though with long
intervals, and Costanza once spent some time in no small disquiet when, during
a severe and lengthy illness, she had had to let her husband write for her,
and no answer came from Brøchner. She then sent him everlastings and red
carnations, which last, according to a remark of Brøchner's, have ``a rich
meaning'' in the Italian language of flowers. In his last years he not
infrequently told his small children anecdotes about \emph{Costanza}, whose
picture they knew. It was hardly without a thought of this name of hers that in
Vilhelm Møller's ``Nyt dansk Maanedsskrift'' [New Danish Monthly] he signed a
treatise with the pseudonym Constans.

\sectionnum{III.}

The ordinary reader finds all too little of Brøchner's nature as a human being
in his writings. The content is too abstract, and the style deters.

By a hair's breadth he would never have become an author at all. He felt no
calling to it; he had let his self-
% --- p. 243 ---
\opage{243}criticism consume his urge to produce without resistance, and by the age of
forty had reached such a point of distaste for all production that he was
glad, or believed he was glad, when elsewhere, in Germany for example, thoughts
were put forward that he had thought and written down for his own use long
before, but had not given to the press. What does it matter, he would say in
such cases, whether it is said by me or by someone else, so long as it comes
out.

At 23 he had published a polemical occasional piece on baptism; only a full
14 years later, to qualify for appointment as docent in philosophy at the
University, did he write, again as an occasional piece, his little book on
Spinoza. Without the persistent pressing of a younger friend it would have
stopped there. But now there followed --- prompted at first by the polemic
against R. Nielsen that others had begun --- after 11 years of silence, at
short intervals, five large volumes and several smaller works, a whole
production written after the author's forty-eighth year, at a time when he was
weakened to the utmost by the pulmonary consumption that for ten years broke
him down inch by inch, and broke him, as in the aggravated death penalty, limb
by limb.

But this very illness, which gave
% --- p. 244 ---
\opage{244}no quarter, lent support to the urging from outside. As long as Brøchner was
still well, or at least did not believe himself mortally stricken, he had many
a time longed for death. Active as he was, he not only had no fear of death
but a kind of love for it. He could speak of it with warmth. It meant for him
liberation from personality, reunion with that impersonal principle of the
life of the All from which the individual is an emanation; and he was fully
content with his self-transformation into a moment in the Godhead's knowledge
of itself in its eternal realm of spirit --- as he mystically put it.

When death approached with rapid steps, it did not thereby take on any
frightening character for him in and of itself; but he strove to live as long
as possible, partly so as to lay the foundation of his children's upbringing in
the spirit in which he wished them brought up, partly so as to leave behind, by
publishing his manuscripts, a testimony of what he had thought and
investigated. So his last years became his only literarily productive ones.

But these large books of his reflect only a part of his human nature, and not
the most human part. They give an idea of his life of thought and of his
pantheistic religiosity. They reveal his
% --- p. 245 ---
\opage{245}acuity and his pathos --- that is all. And, as I said, the style deters. It is
a purely scholastic style, a bookish language with a multitude of German
participial constructions, in which everything is \emph{orderly}, but nothing
has the spontaneity of life. It reflects all that was stiff, exacting and
anxiously punctilious in Brøchner's nature, and gives some idea of a certain
pettiness in small things that in him was not ruled out by the breadth of
vision of his fundamental outlook. But there is a quite different style in his
letters and in his few last popular pieces. Over the stiff style of the large
books there lies a tinge of pedantry that never lay over Brøchner's person. One
sees well enough that he was a strict metaphysician, but one does not detect
his attitude to conventional morality, his free and broad tolerance, in which
something of the man of the world's way of seeing was joined with something of
Spinoza's lofty indifference. One cannot sound the depth or measure the
strength of his loathing for all theology and all theological views when faced
with this official lecture-hall form of his thought, so scientific in its cut
and so dispassionate in its expression. Nor does one see how irritable, how
ready for a sudden break he was with those companions to whom no deeper
sympathy bound him.
% --- p. 246 ---
\opage{246}And who, finally, could infer from these works what Brøchner was as a friend!
If he was tender and lastingly devoted as a lover, and exemplary as a father
who had to take a mother's place with small children, he was in a still higher
degree rare as a friend. He, who was reserved to the point of shyness, and in
conversation always mocking toward any outburst of youthful enthusiasm, was
passionate in friendship. Toward much younger friends he was at once a friend
and a father. He thought unceasingly, year in, year out, with a care that
passed over no detail, about what he could accomplish for his friend. His
friendship was active, a constant activity. And, what is remarkable, it grew
steadily in strength from year to year. It rose in inwardness and
understanding love and reached --- perhaps after fourteen or fifteen years
--- its greatest height in the year he died.

But this is connected with the fundamental trait, which can likewise be made
out only with difficulty through the large books, that Brøchner was in vigorous
development almost until the end came.

It did not look so. He had long been a solitary man. He always knew only a
few people, and in his last years he had
% --- p. 247 ---
\opage{247}withdrawn more and more from the world; nor was he one of those who are much
influenced by their surroundings. He had also, relatively early, taken on
something of the older man's manner. His impressions of life lay fixed in a
stock of anecdotes, to which he often returned without remembering that he had
told them before. He had never been a supple mind, and his antipathies,
literary and scholarly as well as personal, were of an unshakeable nature. His
contempt for a thinker like Schopenhauer went beyond all reasonable measure;
his loathing for the form of national self-deification he found in Grundtvig,
and for the mixture of dilettantism and unreliability he believed he found in
Goldschmidt, knew no bounds. Already in the sixties, when one saw him busy in
his garden in Søgade, intent on developing ever new varieties of pansies, lost
% the print sets "i" twice here, at a line end and again at the start of the next line (transcription.tex, PRINTED AS IS)
in joy at their many different colour patterns, he could give the impression
of a man who was finished with human life and had taken refuge from it in a
life together with plants.

But that was a semblance. In his later years he underwent a marked
metamorphosis. The fighting spirit of his early youth had awakened in him
anew. His reserve had lessened, and his political and social views
% --- p. 248 ---
\opage{248}took on a new colouring. He withdrew completely from people to whom he had
earlier been close in political sympathies; between him and his old party
comrades, the National Liberals, the most decisive break took place in his
mind; he stood far closer to the Socialists than to them.

He had never been inclined to overrate the state of things in Denmark; but in
these last years, when his soul felt the urge to stir its wings, he felt more
strongly the pressure of smallness and pettiness. In a letter from 1870 he
says: ``What there is to tell you from here at home is so little and so poor
that it is hardly worth mentioning. Our public life goes its trivial
parliamentary way, with little deed and much talk, preferably about what is
not worth talking about; with small ideas and small personalities, but with
great illusions, great vanity and great pettinesses. It is not gratifying, but
it must presumably be so, as a passage to a truly free life.''

In a letter from 1874 it reads still more painfully, and as if with a
melancholy look back: ``It is physically unhealthy to live in a small, stuffy
room; it is spiritually unhealthy to live in such dwarfish conditions as here
at home. Everything shrinks so pitifully; one never receives the fresh
impression of what is great; one does not feel
% --- p. 249 ---
\opage{249}the powerful breath of the life of the world's thought. How easily one sinks
down here at home! how easily one takes as the measure of one's striving the
pettiness that surrounds one! how easily one stops in the race when one
notices that one is some way ahead of the few who were competing for the
prize!''

In these years it finally dawned on him that by another manner of writing he
could win himself a wider circle of readers. At times he even went so far as to
regret that as a young man he had not stayed long enough in France to take in
the influence of the French art of exposition --- a remark that could not but
astonish in his mouth, since until then he had had no feeling at all himself
for the shortcomings of the form of his books.

Until then he had, as it were, not dared to write with his whole human nature.
His pamphlet against Nielsen was the first symptom of his urge to strike a more
natural tone. Yet in it the wittily sharp does not quite fuse with the solemn.
In the little treatise on Giordano Bruno there is a clear effort to paint,
strong beginnings of a wholly new manner of narration; and in the posthumous
\textit{Erindringer om S. Kierkegaard} [Recollections of S. Kierkegaard] the
presentation, in all its unpretentiousness, is superbly sharp in its drawing.
% --- p. 250 ---
\opage{250}Every word here renders reality; even the most insignificant thing has life.

He needed only to use the unforced style that he had always had at his
disposal in conversation and in letters, but which he had not dared to trust
himself to.

I had directed a number of objections against his main work, \textit{Bidrag
til Opfattelsen af Filosofiens historiske Udvikling} [Contributions to the
Understanding of the Historical Development of Philosophy]. I had pointed out
the one-sidedness of a presentation that gave only the longitudinal connection
of philosophy, and the inconsistency of Brøchner's sometimes taking the lateral
connection into account and sometimes not. When the subject is Greek
philosophy, he is willing to see it in its connection with the whole Greek
spirit and culture, and holds that only so can it be understood. But when the
subject is modern French, German or English philosophy, it is treated as though
it led an isolated existence, without relation to the national spirit, without
relation to the country's history, poetry and so on. I expressed my admiration
for his gift for seizing the main points and presenting the main features of a
system; but I asked: Is the system, in all philosophers, what is essential and
most important? Is the skeleton everything, the body nothing? Should one not
at least beware of stamping interesting, dissimilar philosophical phenomena
with
% --- p. 251 ---
\opage{251}the same official school-word, when quite different national spirits have
found a voice through them? Condillac is usually presented as an appendix to
Locke; is there not all the same a yawning gulf between their whole ways of
philosophizing? Is there not, on the side of nationality, a still stronger
kinship between Descartes and Condillac than there is, as regards the results of
their thinking, between Condillac and Locke?

In short, I did not hide that the treatment of the history of philosophy that
I had in view as the satisfactory one was a far more concrete and living one
than Brøchner's. He answered by pointing to the slowness with which nationality
among the more recent peoples has worked its way out of the common world of
Latin, and to the extremely complex conditions among the moderns. But if only
his book had been written in the style and with the life that shows itself in
the following passage from a letter:

``Greek philosophy was genuinely national, and therefore could not put forth
\emph{fruitful} shoots in foreign soil. The newer philosophy unfolds on a
larger terrain; like the Indian fig tree it throws out branches on all sides,
which bend to the ground, strike new roots and grow with luxuriant growth.
That is why the complication becomes so great, the demonstration of the
development so involved,
% --- p. 252 ---
\opage{252}that the great features I wanted to bring out easily become indistinct.
% the print has a semicolon after "frem" (have frem; let blive), where the sense wants a comma
Try once to analyse Spinoza from the standpoint of nationality. Try to
determine the mixture of the Jewish, the southern European, the Dutch and the
universally modern. Try to analyse Kant, with his spiritual kinship with the
Germans, his kinship by blood with the Scots, his elective affinity with the
French. You will see the task is not easy; I fear it will prove as difficult as
calculating the mutual perturbations of three heavenly bodies.''

I believe none the less that this was the task that needed solving, and it
seems to me that the very way in which Brøchner wrestles with it here shows
that, with a somewhat different development of his powers, he would have been
able to lay stones for the building he describes as impossible to build.

\sectionnum{IV.}

About a week before the day of his death (he died on 17 December 1875)
Brøchner took leave of his friends. When one day I stepped up to his sickbed, he
put a thick packet into my hand --- it was all the letters I had written to him
over a long run of years, arranged in order of time
% --- p. 253 ---
\opage{253}and marked on the back of the envelope with the date of sending, the date of
receipt and the date on which the letter had been answered. His strong sense of
order, and his care that no confidential communication from a friend should
fall into strangers' hands, show themselves in this small trait.

He had a short conversation with everyone who came, and pressed his hand,
briefly and firmly, in a final farewell.

This last illness, which closed a history of suffering that stretched over
more than a decade, was extremely agonizing.

As his thoughts were drawn toward his approaching death, the sick man came to
occupy himself with the idea of his funeral. The thought of the Church, to
which he had had no relation, and of its priests, who had always been against
him, was very painful to him. And in his feeling of having belonged wholly to
science and of having had no contact with the clergy, he formulated this wish,
which he had laid before the proper authority: to be buried from the vestibule
of the University and not from the Nørrebro chapel. He then expressed the wish
that a younger friend should speak over him instead of the priest, and named
the one he had chosen for it, who had also been willing.

Naturally, after his death there came a
% --- p. 254 ---
\opage{254}refusal with regard to burial from the University's Aula, and with it any
thought of that friend's speaking on the occasion fell away.

Although Brøchner had the liveliest aversion to a priest's speaking over his
coffin --- he said on his deathbed: I believe that if such a black man begins to
preach over me, I shall rise up out of the coffin and protest --- the clerical
address was not avoided.

He was buried on 23 December. The clergyman had, it seems, been given the
impression that Brøchner had wished for his participation; he intimated as
much, and among other things expressed his joy that the deceased and he,
though separated by their convictions, had met in mutual respect. The next day
a then much-read newspaper announced with aplomb that the priest had spoken at
Brøchner's express wish.

For the rest, after his death he was portrayed with precisely the goodwill and
warmth needed to show how great the distance was between an \emph{idealist}
like him and other, younger freethinkers of a quite different mould. Indeed,
the very matter of whether he could rightly be called a freethinker has
gradually become an ever greater question.

A question it is not. But it would testify to narrowness of
% --- p. 255 ---
\opage{255}mind if one held Brøchner's personality to be exhausted in that barren word.
Our opinions are not the most essential or the most original part of our being.
Brøchner was by temperament a searcher, whose passionate inquiry was set in
motion and driven by a \emph{will of spring} like the one that plays its game
in nature when the sap rises in all the plants. As a mind he was a brooder and
a man of natural piety, whose faith in the human spirit and whose relation to
the spirit of the universe had the character of a religion. He stands in the
history of Danish literature as the herald of a new age.

\closingrule

\end{document}
